\section{Introduction}
Local coordination and linker substitution give metal--organic frameworks (MOFs) a chemically diverse electrostatic landscape. Partial atomic charges translate that landscape into parameters used in atomistic simulations, but the assigned values depend on how the electron density is represented. This dependence matters in practice: systematic adsorption calculations show that changing framework point charges can alter carbon dioxide and water isotherms, with particularly strong sensitivity at open metal sites.\cite{Sladekova2020,SladekovaCorrection2021} Understanding how charge definitions respond to local chemistry is therefore useful even when no single definition can serve as a universal reference.

Bader analysis partitions electron density into atomic basins, DDEC6 balances electrostatic and chemical objectives, and CM5 maps Hirshfeld populations through a parameterized charge model.\cite{Tang2009,Manz2016,ManzCorrection2022,Marenich2012} Their differences reflect distinct definitions of an atomic contribution to a continuous density. A signed difference between two assignments consequently describes their relationship; it does not measure an observable charge error. Local environments may organize that relationship by producing directional shifts, larger departures from an average mapping, or a broader range of departures within the same class. These possibilities require separate summaries.

The Quantum MOF (QMOF) Database provides aligned structures and charge assignments for examining these questions at scale.\cite{Rosen2021} Rosen and co-workers already established substantial differences among DDEC6, Bader, and CM5 charges, often exceeding changes associated with the tested density functional approximations.\cite{Rosen2022} Structural charge predictors, including PACMAN, further demonstrate that a specified charge definition can be learned from crystal environments.\cite{Zhao2024,Zhao2024Correction} Neither observation by itself establishes how much chemical information remains in a difference between charge definitions once the central element and one of the charges are known. A structural model can otherwise appear successful by recovering element offsets or curvature between charge scales.

We address this conditional question by comparing matched predictors of an ordered charge difference. A flexible predictor supplied with the element and source charge provides the control; first-shell descriptors and then two-hop path composition provide successive additions. Withholding complete structures tests prediction beyond the fitting records. A second split keeps records with identical source MOFkeys or reduced formulas together, addressing a specified form of related-record dependence. This design makes the chemical contribution assessable without treating prediction error against one charge definition as physical accuracy.

For a common cohort of 9,171 QMOF structures, we combine these comparisons with a paired atlas of chemically interpretable environments. The atlas separates signed departure, absolute departure, and within-environment residual spread. Carbon attachment classes illustrate why these quantities lead to different conclusions, while population-specific evaluation reveals a central-metal failure hidden by the overall improvement. Finite graph tests then delimit the stability of the tested contrasts and identify shared-proton ambiguity. The resulting contribution is a quantitative account of how local chemistry organizes charge-definition disagreement within the evaluated cohort and models.
